% ============================================================================
% SECOND FULL MANUSCRIPT DRAFT
% ============================================================================
% This file is intentionally self-contained so it can be pasted into a LaTeX
% editor and compiled without a separate BibTeX file.
%
% AUTHOR ACTION markers identify information that must be confirmed from the
% parent protocol, the principal investigator, or the final analysis plan.
% The corresponding \draftnote{} boxes remain visible in the compiled document.
%
% Draft prepared from the current MATLAB implementation and the results in:
%   TFA Results New Latest Boss 1/tfa_results.xlsx
% Current analysis cohort: 23 NC and 52 MCI participants.
% ============================================================================

\documentclass[12pt]{article}

\usepackage[T1]{fontenc}
\usepackage[utf8]{inputenc}
\usepackage[margin=1in]{geometry}
\usepackage{amsmath,amssymb}
\usepackage{graphicx}
\usepackage{booktabs}
\usepackage{longtable}
\usepackage{tabularx}
\usepackage{array}
\usepackage{multirow}
\usepackage{pdflscape}
\usepackage{xcolor}
\usepackage{hyperref}
\usepackage{enumitem}
\usepackage{setspace}
\usepackage{caption}
\usepackage{float}

\hypersetup{
    colorlinks=true,
    linkcolor=blue!55!black,
    citecolor=blue!55!black,
    urlcolor=blue!55!black
}

\setstretch{1.35}
\setlength{\parindent}{0.3in}
\setlength{\parskip}{0.15em}
\setcounter{secnumdepth}{3}
\setcounter{tocdepth}{2}

\newcolumntype{Y}{>{\raggedright\arraybackslash}X}
\newcommand{\PETCOtwo}{P_{\mathrm{ET}}CO_2}
\newcommand{\CBFV}{\mathrm{CBFV}}
\newcommand{\MAP}{\mathrm{MAP}}
\newcommand{\COtwo}{CO_2}
\newcommand{\Real}{\operatorname{Re}}
\newcommand{\Imag}{\operatorname{Im}}
\newcommand{\atanTwo}{\operatorname{atan2}}

% A visible box for decisions that should not be silently filled in.
\newcommand{\draftnote}[1]{%
    \par\medskip
    \noindent\fcolorbox{orange!70!black}{orange!8}{%
        \parbox{0.94\linewidth}{\textbf{DRAFT NOTE:} #1}}
    \par\medskip
}

% A visible figure placeholder. The optional first argument controls height.
\newcommand{\figureplaceholder}[4][0.22\textheight]{%
    \begin{figure}[p]
        \centering
        \fcolorbox{gray!80}{gray!6}{%
            \parbox[c][#1][c]{0.91\linewidth}{%
                \centering
                \textbf{FIGURE PLACEHOLDER: #2}\par\medskip
                #3}}
        \caption{#2. Draft placeholder; replace this caption after the final
        figure is prepared.}
        \label{#4}
    \end{figure}
}

\title{Separating Mean Arterial Pressure- and Carbon Dioxide-Associated
Contributions to Cerebral Blood Flow Velocity in Mild Cognitive Impairment
Using Multivariable Transfer Function Analysis}

% AUTHOR ACTION: Replace placeholders with the final author list and affiliations.
\author{[Author names and affiliations to be added]}
\date{Second full draft -- \today}

\begin{document}

\maketitle

\draftnote{Confirm the final title, author order, affiliations, corresponding
author, target journal, and whether the journal prefers ``cerebral blood
velocity'' or ``cerebral blood flow velocity.'' This draft uses
``associated'' instead of ``mediated'' because the observational transfer
functions do not establish causality.}

\begin{abstract}
\textbf{Background:} Dynamic cerebral autoregulation is commonly assessed
with single-input, single-output (SISO) transfer function analysis using mean
arterial pressure (MAP) as the input and cerebral blood flow velocity (CBFV)
as the output. CBFV is also affected by carbon dioxide, so a MAP-only model
may mix pressure-associated and carbon-dioxide-associated dynamics.

\textbf{Objective:} We tested whether simultaneously modeling MAP and
end-tidal carbon dioxide ($\PETCOtwo$) materially changed pathway-specific
transfer function estimates and whether the resulting multivariable estimates
differed between cognitively normal controls (NC) and participants with mild
cognitive impairment (MCI).

\textbf{Methods:} Baseline MAP, $\PETCOtwo$, and CBFV recordings from 23 NC
and 52 MCI participants were resampled to 4~Hz. Auto- and cross-spectra were
estimated using 128-second Hann windows with 50\% overlap and at least three
Welch windows. Separate SISO models and a two-input, single-output (MISO)
model were calculated over 0--0.35~Hz. Prespecified bands were 0.005--0.024,
0.024--0.070, 0.070--0.200, and 0.200--0.350~Hz. Gain comparisons used paired
tests within participants; NC--MCI comparisons used unequal-variance tests.
Phase was analyzed with circular statistics and permutation tests. The
Benjamini--Hochberg (BH) method controlled the false discovery rate within
prespecified analysis families.

\textbf{Results:} Simultaneous modeling changed gain estimates in both
groups. MISO MAP gain was lower than SISO MAP gain in the very-very-low-frequency
band in NC (mean difference $-0.240$, 95\% confidence interval [CI]
$-0.384$ to $-0.096$, BH-adjusted $P=0.0052$) and MCI ($-0.183$, 95\% CI
$-0.290$ to $-0.076$, BH-adjusted $P=0.0038$). The larger and more consistent
changes occurred in the carbon dioxide pathway. MISO $\COtwo$ gain was lower
than SISO $\COtwo$ gain in the low- and high-frequency bands in NC and in the
very-low-, low-, and high-frequency bands in MCI (all BH-adjusted $P<0.006$).
No prespecified band-level NC--MCI difference in MISO gain, partial coherence,
phase, or MAP-minus-$\COtwo$ pathway balance survived BH correction. An
exploratory frequency-wise analysis identified higher NC than MCI MISO MAP
gain at 0.1484~Hz (BH-adjusted $P=0.0366$ within that frequency curve).

\textbf{Conclusions:} Modeling MAP and $\PETCOtwo$ together materially
changed pathway-specific gain estimates, particularly for the $\COtwo$
pathway. These results support the methodological premise that MAP-only and
$\COtwo$-only SISO estimates are not interchangeable with conditional MISO
estimates. However, the present data did not support the prespecified
hypothesis of a broad band-level difference in conditional cerebral
hemodynamics between NC and MCI. The isolated 0.1484-Hz group difference is
hypothesis-generating and requires confirmation.
\end{abstract}

\noindent\textbf{Keywords:} cerebral autoregulation; carbon dioxide; mild
cognitive impairment; transcranial Doppler; transfer function analysis;
multivariable spectral analysis; coherence

\newpage
\tableofcontents
\newpage

% ============================================================================
\section{Introduction}
% ============================================================================

The brain requires a continuous supply of oxygen and nutrients, but systemic
blood pressure changes from moment to moment. Cerebral autoregulation is one
part of the broader set of mechanisms that helps maintain cerebral perfusion
despite these changes \cite{Claassen2021}. Dynamic cerebral autoregulation
(dCA) focuses on the response of cerebral blood flow to relatively rapid or
spontaneous changes in arterial pressure. In human studies, dCA is often
estimated from continuous MAP and transcranial Doppler-derived CBFV signals.
Transfer function analysis (TFA) provides a frequency-specific description of
this pressure--flow relationship \cite{Panerai2023}.

Cerebrovascular regulation is relevant to mild cognitive impairment (MCI)
because vascular dysfunction and reduced perfusion may contribute to the
development or progression of cognitive decline \cite{Kisler2017,Beishon2017}.
At the same time, the existing dCA literature in MCI is mixed. Individual
studies have often found limited or no group differences in conventional
TFA gain and phase \cite{Gommer2012,Tarumi2014}, and a recent systematic review
did not find consistent evidence of poorer spontaneous dCA in MCI or
Alzheimer's disease \cite{Heutz2023}. These mixed results may reflect true
biological heterogeneity, limited sample sizes, different experimental
conditions, or limitations in how the input--output system is represented.

Standard TFA usually treats MAP as the only input to CBFV. This is useful and
well established, but it is incomplete. Carbon dioxide is a strong and rapid
regulator of cerebral vascular tone, and even small changes in arterial carbon
dioxide can produce substantial changes in cerebral blood flow
\cite{Hoiland2019}. End-tidal carbon dioxide is not identical to arterial
carbon dioxide, but it is commonly used as a practical continuous surrogate.
If MAP and $\PETCOtwo$ fluctuate in overlapping frequency ranges, a MAP-only
SISO transfer function can contain CBFV variation associated with both MAP
and $\COtwo$. The same problem applies in reverse to a $\COtwo$-only SISO
model.

Prior multivariable work provides a basis for addressing this problem.
Panerai and colleagues modeled spontaneous MAP and $\PETCOtwo$ fluctuations
together \cite{Panerai2000}. Peng and colleagues later showed that including
gas-related inputs increased multiple coherence below 0.05~Hz and could alter
the low-frequency MAP-to-CBFV transfer function \cite{Peng2008}. Their result
does not mean that every multivariable estimate is automatically more
accurate. It shows that the univariable estimate and the conditional
multivariable estimate answer different questions. The SISO model describes
the total pairwise association between one input and CBFV. The MISO model
estimates the association of each input with CBFV while accounting for the
other measured input.

This distinction is important for MCI. A MAP-only group comparison may miss
or misstate a pressure-associated difference if the groups also differ in
$\COtwo$ dynamics or in MAP--$\COtwo$ coupling. Conversely, a multivariable
model may show that an apparent SISO difference is shared with the second
input. The purpose of the present study was therefore not simply to maximize
coherence. It was to determine whether simultaneous MAP and $\PETCOtwo$
modeling materially changes pathway-specific interpretation, and then to use
that model to compare NC and MCI participants.

\subsection{Central hypothesis and aims}

Our central hypothesis was that simultaneously modeling MAP and $\PETCOtwo$
would materially alter pathway-specific transfer function estimates relative
to separate SISO models, particularly where the two inputs share frequency
content, and that these conditional estimates could reveal whether MAP- and
$\COtwo$-associated CBFV dynamics differ between NC and MCI.

The aims were distinct tests of this hypothesis:

\begin{enumerate}[leftmargin=0.45in]
    \item \textbf{Model dependence and pathway separation.} We compared
    MISO and SISO gain and phase within the same participants. We expected
    the largest changes where one-input estimates contained information
    shared with the second input.

    \item \textbf{Input coupling.} We tested whether MAP--$\COtwo$ input
    coherence was associated with the absolute size of the MISO-minus-SISO
    gain change. This analysis was methodological and exploratory because
    the two quantities are derived from the same spectral matrix.

    \item \textbf{Main biological comparison.} We compared NC and MCI using
    MISO MAP gain, MAP partial coherence, and circular phase after accounting
    for $\PETCOtwo$. Parallel $\COtwo$-pathway comparisons tested whether any
    group pattern was pressure-specific.

    \item \textbf{Pathway balance.} We compared the within-subject difference
    between MAP and $\COtwo$ partial coherence across groups. Positive values
    indicate a stronger conditional MAP--CBFV association than conditional
    $\COtwo$--CBFV association in that band.
\end{enumerate}

Multiple coherence was retained as a descriptive measure of how much CBFV
power was linearly represented in the fitted two-input model. It was not
treated as proof that MISO predicts unseen data better than SISO. Such a claim
would require cross-validation on recordings long enough to estimate a model
and evaluate it on an independent segment.

% AUTHOR ACTION: Create a clean schematic, not a screenshot from MATLAB.
\figureplaceholder[0.22\textheight]
{Conceptual difference between SISO and MISO transfer function analysis}
{Show MAP-to-CBFV and $\PETCOtwo$-to-CBFV as two separate SISO models on the
left. On the right, show MAP and $\PETCOtwo$ entering one MISO model, with a
dashed line between the inputs representing their cross-spectrum. State that
SISO estimates pairwise associations, whereas MISO estimates each pathway
conditional on the other measured input.}
{fig:concept}

% ============================================================================
\section{Methods}
% ============================================================================

\subsection{Study design and participants}

This was a cross-sectional secondary analysis of baseline physiological
recordings from NC and MCI participants. The final TFA cohort contained 23 NC
and 52 MCI participants. All 75 recordings met the analysis requirement of at
least three Welch windows after preprocessing. Demographic information was
available for all NC participants and 51 of 52 MCI participants; MCI 302 was
present in the physiological dataset but did not have a matching row in the
available demographic workbook.

The available cognitive descriptors included the Mini-Mental State
Examination (MMSE) \cite{Folstein1975} and Montreal Cognitive Assessment
(MoCA) \cite{Nasreddine2005}.

% AUTHOR ACTION: Replace this box using the parent protocol and PI-approved text.
\draftnote{Add the parent study name, recruitment setting, exact NC and MCI
diagnostic criteria, inclusion and exclusion criteria, institutional review
board approval number, informed-consent language, and whether diagnosis was
amnestic, non-amnestic, single-domain, or multidomain MCI. Confirm how CDR,
MoCA, neuropsychological testing, and clinical adjudication were used. Do not
infer these details from the spreadsheet alone.}

\subsection{Physiological recording and signal preparation}

Continuous arterial pressure, middle cerebral artery CBFV, end-tidal carbon
dioxide, and electrocardiography were acquired during resting baseline. The
source data were recorded at 500~Hz in the acquisition system and converted
before the present analysis into time-aligned beat-to-beat MAP and CBFV and
breath-to-breath $\PETCOtwo$ values. Cardiac R peaks defined the beat intervals.
The available beat-level signals were visually reviewed before export. No
short artifacts required interpolation in the recordings included here.

% AUTHOR ACTION: Confirm all hardware and acquisition details.
\draftnote{Add the recording posture, baseline duration target, arterial
pressure device and calibration, transcranial Doppler manufacturer and probe
frequency, insonated artery and side, probe fixation method, capnograph model,
sampling hardware/software versions, and room/breathing conditions. Confirm
whether CBFV was unilateral or averaged across hemispheres.}

\subsection{Preprocessing}

The MATLAB workflow first removed rows containing non-finite values in any of
the time, MAP, $\PETCOtwo$, or CBFV signals. Some exported $\PETCOtwo$ columns
contained leading zeros and a large startup transition before stable values
were present. The software removed this startup portion automatically so that
columns did not need to be edited manually. It identified the first nonzero
$\PETCOtwo$ value, calculated absolute consecutive changes in the remaining
signal, and defined unusually large changes using the mean plus one standard
deviation of the upper 10\% of changes. Samples were removed from the start
until a nonzero value followed by a change below this data-derived threshold
was reached.

After startup removal, all signals were linearly resampled to 4~Hz. Mean MAP,
mean $\PETCOtwo$, and mean CBFV were stored before centering. CBFV was then
normalized to percent of its recording mean,

\begin{equation}
    \CBFV_{\%}(t)=100\frac{\CBFV(t)}{\overline{\CBFV}}.
    \label{eq:cbfv_normalization}
\end{equation}

The mean was removed from MAP, $\PETCOtwo$, and normalized CBFV before spectral
estimation. Mean removal eliminates the constant, zero-frequency component so
that the spectra describe fluctuations around each participant's baseline.
Polynomial detrending was available in the software but disabled for the
present analysis to avoid imposing an additional low-frequency transformation.

\subsection{Welch spectral estimation}

Auto- and cross-spectra were estimated using Welch's method
\cite{Welch1967}. Each segment used a 128-second Hann window with 50\% overlap.
At 4~Hz, the window length was

\begin{equation}
    L = 128\ \mathrm{s}\times 4\ \mathrm{samples/s}=512\ \mathrm{samples},
\end{equation}

with an overlap of 256 samples and a step of 256 samples. Therefore, the
minimum number of points required for three windows was

\begin{equation}
    N_{\min}=L+(K-1)(L-L_{\mathrm{overlap}})
    =512+2(256)=1024\ \mathrm{samples},
\end{equation}

which corresponds to 256 seconds at 4~Hz. The fast Fourier transform length
was 512 samples, giving a frequency spacing of

\begin{equation}
    \Delta f=\frac{f_s}{N_{\mathrm{FFT}}}
    =\frac{4}{512}=0.0078125\ \mathrm{Hz}.
\end{equation}

The current full-frequency outputs covered 0--0.35~Hz. With the available
frequency spacing, the last included frequency was 0.34375~Hz. The same
three-point smoothing kernel,

\begin{equation}
    \mathbf{w}=\begin{bmatrix}0.25 & 0.50 & 0.25\end{bmatrix},
    \label{eq:smoothing_kernel}
\end{equation}

was convolved with every auto-spectrum and cross-spectrum. Applying the same
linear smoothing to all spectral quantities preserves their alignment and
reduces isolated bin-to-bin variation. The smoothing setting and kernel were
parameterized in the main analysis file.

For complex spectra, this paper uses the convention

\begin{equation}
    S_{xy}(f)=\mathbb{E}\!\left[X^*(f)Y(f)\right],
    \label{eq:cross_spectrum_convention}
\end{equation}

where $^*$ denotes complex conjugation. This convention matters because the
order of the signals determines the sign of phase.

\subsection{Frequency bands}

Band values were calculated after each subject's frequency-domain model was
estimated. Each subject therefore contributed one value per band to each
band-level statistical test; frequency bins were not treated as independent
participants. Arithmetic metrics were averaged normally across bins. Phase
was averaged circularly as described below. The lower band edge was included
and the upper edge was excluded, except that the final band included its upper
edge.

\begin{table}[H]
\centering
\caption{Prespecified frequency bands used in the analysis.}
\label{tab:bands}
\begin{tabular}{lll}
\toprule
Band & Abbreviation & Frequency range (Hz) \\
\midrule
Very-very-low frequency & VVLF & $0.005 \leq f < 0.024$ \\
Very-low frequency & VLF & $0.024 \leq f < 0.070$ \\
Low frequency & LF & $0.070 \leq f < 0.200$ \\
High frequency & HF & $0.200 \leq f \leq 0.350$ \\
\bottomrule
\end{tabular}
\end{table}

\subsection{Single-input, single-output transfer functions}

For a SISO model with input $x(t)$ and output $y(t)$, the frequency-domain
model is

\begin{equation}
    Y(f)=H_{xy}(f)X(f)+E(f),
    \label{eq:siso_model}
\end{equation}

where $H_{xy}(f)$ is the transfer function and $E(f)$ is unexplained output.
Multiplying both sides by $X^*(f)$ and taking the expectation gives

\begin{align}
    \mathbb{E}[X^*Y]
    &=H_{xy}\mathbb{E}[X^*X]+\mathbb{E}[X^*E],\\
    S_{xy}
    &=H_{xy}S_{xx},
\end{align}

assuming the residual is uncorrelated with the input. The transfer function
is therefore

\begin{equation}
    H_{xy}(f)=\frac{S_{xy}(f)}{S_{xx}(f)}.
    \label{eq:siso_transfer}
\end{equation}

Separate SISO transfer functions were calculated for MAP-to-CBFV and
$\PETCOtwo$-to-CBFV:

\begin{equation}
    H_{\MAP,y}^{\mathrm{SISO}}
    =\frac{S_{\MAP,y}}{S_{\MAP,\MAP}},
    \qquad
    H_{\COtwo,y}^{\mathrm{SISO}}
    =\frac{S_{\COtwo,y}}{S_{\COtwo,\COtwo}}.
\end{equation}

Gain was the magnitude of the complex transfer function,

\begin{equation}
    G_{xy}(f)=\left|H_{xy}(f)\right|.
    \label{eq:gain}
\end{equation}

Because CBFV was expressed as percent of its recording mean and both inputs
were expressed in mmHg, MAP and $\COtwo$ gains are reported as percentage
points of baseline CBFV per mmHg.

Wrapped phase was calculated explicitly as

\begin{equation}
    \phi_{xy}(f)=\atanTwo\!\left(
    \Imag\{H_{xy}(f)\},
    \Real\{H_{xy}(f)\}\right),
    \label{eq:phase}
\end{equation}

where $\Imag\{\cdot\}$ denotes the imaginary part. In standard notation,
Equation~\ref{eq:phase} is the complex angle of $H$ and
returns a wrapped value in $[-\pi,\pi]$. Pairwise magnitude-squared coherence
was

\begin{equation}
    \gamma^2_{xy}(f)=
    \frac{|S_{xy}(f)|^2}{S_{xx}(f)S_{yy}(f)}.
    \label{eq:ordinary_coherence}
\end{equation}

The unexplained fraction was $1-\gamma^2_{xy}$, and the corresponding residual
power was

\begin{equation}
    S_{ee}(f)=S_{yy}(f)\left[1-\gamma^2_{xy}(f)\right].
\end{equation}

The CARNet ordinary-coherence reference for three Welch windows is 0.51 and
for four windows is 0.41 \cite{Panerai2023}. These values were plotted only as
visual benchmarks on ordinary SISO coherence figures. No frequency bin or
participant was filtered on the basis of coherence. The ordinary-coherence
reference was not applied to partial or multiple coherence because those
statistics have different sampling behavior.

\subsection{Two-input, single-output transfer functions}

The MISO model treated MAP as $x_1(t)$, $\PETCOtwo$ as $x_2(t)$, and normalized
CBFV as $y(t)$:

\begin{equation}
    Y(f)=H_1(f)X_1(f)+H_2(f)X_2(f)+E(f).
    \label{eq:miso_model}
\end{equation}

Multiplication by each conjugated input and expectation produces the spectral
normal equations

\begin{align}
    S_{1y} &= S_{11}H_1+S_{12}H_2,\\
    S_{2y} &= S_{21}H_1+S_{22}H_2.
\end{align}

In matrix form,

\begin{equation}
    \underbrace{
    \begin{bmatrix}
        S_{11}(f) & S_{12}(f)\\
        S_{21}(f) & S_{22}(f)
    \end{bmatrix}}_{\mathbf{S}_{xx}(f)}
    \underbrace{
    \begin{bmatrix}
        H_1(f)\\H_2(f)
    \end{bmatrix}}_{\mathbf{H}(f)}
    =
    \underbrace{
    \begin{bmatrix}
        S_{1y}(f)\\S_{2y}(f)
    \end{bmatrix}}_{\mathbf{S}_{xy}(f)}.
    \label{eq:miso_matrix}
\end{equation}

If the determinant

\begin{equation}
    D(f)=S_{11}(f)S_{22}(f)-S_{12}(f)S_{21}(f)
\end{equation}

is nonzero, the two transfer functions can be written as

\begin{align}
    H_1(f)&=\frac{S_{22}S_{1y}-S_{12}S_{2y}}{D},
    \label{eq:h1_explicit}\\
    H_2(f)&=\frac{S_{11}S_{2y}-S_{21}S_{1y}}{D}.
    \label{eq:h2_explicit}
\end{align}

The implementation did not calculate the inverse or Equations
\ref{eq:h1_explicit}--\ref{eq:h2_explicit} directly. At each frequency, MATLAB
solved

\begin{equation}
    \mathbf{H}(f)=\mathbf{S}_{xx}(f)\backslash\mathbf{S}_{xy}(f),
    \label{eq:miso_solve}
\end{equation}

using matrix left division. No ridge regularization or special singular-matrix
fallback was applied in the present analysis. Gain and phase for $H_1$ and
$H_2$ were calculated using Equations~\ref{eq:gain} and \ref{eq:phase}.

\subsubsection{Multiple coherence}

The fitted output power represented by both inputs is

\begin{equation}
    S_{\widehat{y}\widehat{y}}(f)
    =\mathbf{S}_{xy}^{H}(f)\mathbf{H}(f),
\end{equation}

where the superscript $H$ denotes the conjugate transpose. Multiple coherence
was calculated as

\begin{equation}
    \gamma^2_{y\cdot 12}(f)=
    \frac{\Real\left\{\mathbf{S}_{xy}^{H}(f)
    \mathbf{H}(f)\right\}}{S_{yy}(f)}.
    \label{eq:multiple_coherence}
\end{equation}

The MISO unexplained fraction was $1-\gamma^2_{y\cdot 12}$ and residual power
was $S_{yy}(1-\gamma^2_{y\cdot 12})$. Because Equation
\ref{eq:multiple_coherence} is calculated on the same data used to estimate
the model, it was interpreted as in-sample linear representation, not
out-of-sample predictive accuracy.

\subsubsection{Partial coherence}

Partial coherence measures the remaining linear association between one input
and the output after removing the linear component associated with the other
input. For MAP and CBFV conditional on $\PETCOtwo$, the residual cross-spectrum
and auto-spectra were

\begin{align}
    S_{1y\cdot2}
    &=S_{1y}-\frac{S_{12}S_{2y}}{S_{22}},\\
    S_{11\cdot2}
    &=S_{11}-\frac{S_{12}S_{21}}{S_{22}},\\
    S_{yy\cdot2}
    &=S_{yy}-\frac{S_{y2}S_{2y}}{S_{22}}.
\end{align}

The conditional MAP--CBFV partial coherence was then

\begin{equation}
    \gamma^2_{1y\mid2}(f)=
    \frac{|S_{1y\cdot2}(f)|^2}
    {\Real\{S_{11\cdot2}(f)S_{yy\cdot2}(f)\}}.
    \label{eq:partial_map}
\end{equation}

The conditional $\COtwo$--CBFV partial coherence, $\gamma^2_{2y\mid1}$, was
calculated by reversing the roles of inputs 1 and 2. Partial coherence was not
described as the fraction of CBFV caused by a pathway. It is a conditional
linear association after accounting for the other measured input.

\subsubsection{Input coherence and conditioning}

MAP--$\PETCOtwo$ input coherence was

\begin{equation}
    \gamma^2_{12}(f)=
    \frac{|S_{12}(f)|^2}{S_{11}(f)S_{22}(f)}.
    \label{eq:input_coherence}
\end{equation}

The numerical condition number of the raw input spectral matrix was also
stored:

\begin{equation}
    \kappa\!\left(\mathbf{S}_{xx}(f)\right)
    =\|\mathbf{S}_{xx}(f)\|\,
    \|\mathbf{S}_{xx}^{-1}(f)\|.
    \label{eq:condition_number}
\end{equation}

A large value indicates that the transfer function solution can be sensitive
to small spectral changes. However, the raw matrix contains MAP and
$\PETCOtwo$ in different units and power scales, so its condition number is
also scale-dependent. It was used as a diagnostic, not an exclusion rule or a
direct measure of physiological collinearity.

\subsection{Phase representation and circular statistics}

Wrapped phase was the primary representation for statistical analysis.
Standard MATLAB unwrapping was used for the displayed unwrapped curves. The
algorithm adds or subtracts $2\pi$ when it detects a discontinuity between
neighboring frequency bins. This changes the displayed branch of the curve
without changing its circular direction. An experimental custom unwrapping
method remained available in the software but was not used for the reported
results.

Linear averaging is inappropriate for wrapped angles because values just
below $\pi$ and just above $-\pi$ are close on a circle but far apart
arithmetically. For phase values $\theta_1,\ldots,\theta_n$, the mean cosine
and sine components were

\begin{equation}
    C=\frac{1}{n}\sum_{j=1}^{n}\cos\theta_j,
    \qquad
    S=\frac{1}{n}\sum_{j=1}^{n}\sin\theta_j.
\end{equation}

The circular mean direction was

\begin{equation}
    \overline{\theta}_{\mathrm{circ}}=\atanTwo(S,C),
    \label{eq:circular_mean}
\end{equation}

and the mean resultant length and circular standard deviation were

\begin{equation}
    R=\sqrt{C^2+S^2},
    \qquad
    s_{\mathrm{circ}}=\sqrt{-2\ln R}.
    \label{eq:circular_sd}
\end{equation}

These are standard directional-statistics definitions
\cite{MardiaJupp2000,Berens2009}. Group phase means were first calculated
circularly from the wrapped subject phases at each frequency. The circular
mean curve was then unwrapped for display. The same circular standard
deviation was retained around the wrapped and unwrapped group means because
adding complete phase cycles does not change circular spread.

\subsection{Statistical analysis}

Statistical analysis was performed at two levels. Prespecified inference used
one value per subject per band. Frequency-wise tests were exploratory and were
performed separately at each frequency bin.

All inferential tests were two-sided. Confidence intervals were 95\%, and the
nominal significance level was 0.05 before the stated multiple-comparison
adjustments.

\subsubsection{Participant characteristics}

Continuous participant characteristics were summarized as mean, standard
deviation, and range and compared using Welch's unequal-variance two-sample
$t$ test \cite{Welch1947}. Sex and medication counts were compared with
Fisher's exact test. These descriptive comparisons were not treated as the
primary TFA hypothesis family.

\subsubsection{MISO versus SISO comparisons}

MISO and SISO values came from the same subjects. Arithmetic outcomes were
therefore compared using paired $t$ tests on

\begin{equation}
    d_i=x_i^{\mathrm{MISO}}-x_i^{\mathrm{SISO}}.
\end{equation}

The null hypothesis was $H_0:\mu_d=0$. The standardized paired effect was

\begin{equation}
    d_z=\frac{\overline{d}}{s_d}.
\end{equation}

Because gain ratios were skewed in some bands, the median subject-level ratio

\begin{equation}
    \widetilde{R}=\operatorname{median}_i
    \left(\frac{x_i^{\mathrm{MISO}}}{x_i^{\mathrm{SISO}}}\right)
    \label{eq:median_ratio}
\end{equation}

and its interquartile range were reported as supporting effect descriptions.
Paired phase differences were wrapped to the shortest angular difference and
tested with a sign-randomization circular permutation test.

\subsubsection{NC versus MCI comparisons}

For arithmetic outcomes, NC and MCI were compared with Welch's two-sample
$t$ test. The test statistic was

\begin{equation}
    t=\frac{\overline{x}_{\mathrm{NC}}-
    \overline{x}_{\mathrm{MCI}}}
    {\sqrt{s_{\mathrm{NC}}^2/n_{\mathrm{NC}}+
    s_{\mathrm{MCI}}^2/n_{\mathrm{MCI}}}}.
    \label{eq:welch_t}
\end{equation}

The null hypothesis was equality of group means at the tested band or
frequency. Hedges' $g$ described the standardized group difference with a
small-sample correction \cite{Hedges1981}. Independent circular phase tests
randomly reassigned pooled wrapped phase values to groups while preserving
group sizes.

\subsubsection{Permutation settings}

Circular tests used 10,000 permutations. For each test, the raw permutation
$P$ value was

\begin{equation}
    P=\frac{1+\#\left(T_{\mathrm{perm}}\geq
    T_{\mathrm{observed}}\right)}{10{,}000+1}.
\end{equation}

The added one prevents a zero $P$ value. A fixed random seed of 2026 made the
random permutations reproducible; it did not pull results toward zero or
change the null hypothesis.

\subsubsection{Multiple-comparison control}

Raw $P$ values answer one test at a time. Because many related bands and
metrics were tested, the Benjamini--Hochberg procedure controlled the expected
false discovery rate within prespecified analysis families
\cite{BenjaminiHochberg1995}. For sorted raw values
$P_{(1)}\leq\cdots\leq P_{(m)}$, the adjusted values were calculated as

\begin{equation}
    P^{\mathrm{BH}}_{(i)}=
    \min_{j\geq i}\left\{1,\frac{m}{j}P_{(j)}\right\}.
    \label{eq:bh}
\end{equation}

Statistical significance was defined as $P^{\mathrm{BH}}<0.05$. A value above
0.05 was interpreted as insufficient evidence to reject the null, not as
proof that the groups or models were equivalent.

The primary MISO-versus-SISO family contained 16 gain comparisons: two groups,
two pathways, and four bands. The primary NC--MCI family contained 16 MISO
comparisons: two pathways, two metrics (gain and partial coherence), and four
bands. The 32 secondary within-subject model comparisons (coherence and phase)
were adjusted together. The 36 secondary group comparisons (MISO phase, all
SISO outcomes, and four pathway-balance comparisons) were adjusted together.
The 24 input-coherence associations were one exploratory family. The nine
physiological recording-characteristic tests exported by the workflow were
also adjusted together, although Table~\ref{tab:participants} reports raw
descriptive $P$ values for participant and recording characteristics.

\subsubsection{Exploratory frequency-wise tests}

At each included frequency, NC and MCI were compared for MISO gain, partial
coherence, and wrapped phase. MISO and SISO were also compared within each
group for gain and phase. BH adjustment was applied across the bins belonging
to each individual metric/pathway comparison curve. Thus, the adjusted
$P$-value curve shows where one specific frequency curve provides evidence
against its null hypothesis after accounting for the bins on that curve. The
adjustment did not additionally control across every metric and pathway in
the full project.

% AUTHOR ACTION: This is a genuine analysis-plan decision, not just wording.
\draftnote{The current exploratory frequency-wise analysis includes the
0-Hz bin because the configured analysis range begins at 0~Hz, whereas the
prespecified bands begin at 0.005~Hz. Mean removal should minimize the DC
component, but the final paper should decide in advance whether to retain or
exclude 0~Hz from the exploratory family and then rerun that analysis if the
decision changes.}

\subsection{Software and reproducibility}

The analysis was implemented in MATLAB R2025b using a single workflow for
single-subject, batch, and synthetic-data runs. The same Welch spectra were
passed to enabled SISO and MISO models, and the resulting subject-level values
were passed to plotting, group summarization, statistics, and Excel export.
The source code includes unit tests for preprocessing, spectral estimation,
SISO and MISO calculations, circular statistics, BH adjustment, workflow
settings, statistical analysis, and Excel export.

% AUTHOR ACTION: Add a permanent repository URL and commit identifier at submission.
\draftnote{Add the final public repository URL, release/commit identifier,
MATLAB toolbox requirements, data-availability statement, and instructions
for obtaining any restricted participant data.}

% ============================================================================
\section{Results}
% ============================================================================

\subsection{Participants and recording characteristics}

All 23 NC and 52 MCI recordings completed preprocessing and TFA. The NC and
MCI groups did not differ in usable recording duration, number of Welch
windows, mean MAP, mean $\PETCOtwo$, mean CBFV, or CVRi (Table
\ref{tab:participants}). Demographic matching was available for 23 NC and 51
MCI participants. Age, sex, body size, education, MMSE, blood pressure, heart
rate, and medication counts were similar between groups. As expected from the
clinical grouping, MoCA was lower in MCI. The available geriatric depression
score was also higher in MCI, although values were low in both groups.

\begin{longtable}{p{0.30\linewidth}p{0.25\linewidth}p{0.25\linewidth}p{0.10\linewidth}}
\caption{Participant and recording characteristics.}
\label{tab:participants}\\
\toprule
Characteristic & NC & MCI & Raw $P$ \\
\midrule
\endfirsthead
\toprule
Characteristic & NC & MCI & Raw $P$ \\
\midrule
\endhead
\multicolumn{4}{l}{\textit{Demographic and clinical characteristics}}\\
Age, years & $66.15\pm7.14$ [55.98--80.53] & $64.78\pm6.12$ [55.11--78.75] & 0.430 \\
Female sex, $n/N$ (\%) & 13/23 (56.5) & 29/51 (56.9) & 1.000 \\
Height, cm & $169.83\pm7.98$ [157.5--186.5] & $168.35\pm9.05$ [143.5--187.8] & 0.482 \\
Weight, kg & $78.59\pm16.67$ [50.6--110.6] & $78.90\pm14.71$ [49.9--116.8] & 0.938 \\
Body mass index, kg/m$^2$ & $27.10\pm4.65$ [20.14--35.31] & $27.74\pm4.03$ [20.27--35.81] & 0.573 \\
Education, years & $16.22\pm2.37$ [12--20] & $15.65\pm2.39$ [12--19] & 0.345 \\
MMSE, points & $29.17\pm0.94$ [27--30] & $29.18\pm1.14$ [24--30] & 0.992 \\
MoCA, points & $27.35\pm1.85$ [24--30] & $24.90\pm2.72$ [17--30] & $<0.001$ \\
Systolic BP, mmHg & $129.65\pm10.69$ [113--153] & $131.63\pm11.29$ [106--158] & 0.474 \\
Diastolic BP, mmHg & $73.70\pm8.90$ [61--93] & $74.61\pm7.88$ [56--91] & 0.675 \\
Heart rate, beats/min & $72.35\pm7.37$ [59--86] & $70.51\pm10.04$ [50--87] & 0.381 \\
Blood-pressure medication, $n/N$ & 8/23 & 16/47 & 1.000 \\
Cholesterol medication, $n/N$ & 8/23 & 9/47 & 0.234 \\
Geriatric depression score & $0.57\pm0.95$ [0--3] & $1.94\pm1.78$ [0--7]$^{a}$ & $<0.001$ \\
\addlinespace
\multicolumn{4}{l}{\textit{Physiological recording and analysis characteristics}}\\
Mean MAP, mmHg & $90.73\pm14.47$ [68.64--116.97] & $95.26\pm13.90$ [67.94--128.48] & 0.214 \\
Within-recording MAP SD, mmHg & $4.11\pm1.14$ [2.31--7.19] & $4.17\pm1.44$ [1.78--9.50] & 0.843 \\
Mean $\PETCOtwo$, mmHg & $36.68\pm3.46$ [29.79--42.50] & $36.56\pm2.85$ [31.01--43.04] & 0.890 \\
Within-recording $\PETCOtwo$ SD, mmHg & $1.41\pm0.68$ [0.58--3.17] & $1.32\pm0.58$ [0.48--3.31] & 0.597 \\
Mean CBFV, cm/s & $46.43\pm11.75$ [31.98--81.99] & $46.83\pm10.17$ [30.72--82.26] & 0.888 \\
Within-recording CBFV SD, cm/s & $3.55\pm1.16$ [1.91--6.50] & $3.35\pm1.36$ [1.35--8.66] & 0.519 \\
CVRi, mmHg/(cm/s) & $2.03\pm0.45$ [1.06--2.89] & $2.14\pm0.60$ [1.13--3.76] & 0.413 \\
Usable duration, s & $297.35\pm13.02$ [270.5--347.8] & $297.43\pm14.19$ [260.5--338.8] & 0.980 \\
Welch windows, $n$ & $3.04\pm0.21$ [3--4] & $3.06\pm0.24$ [3--4] & 0.795 \\
\bottomrule
\multicolumn{4}{p{0.94\linewidth}}{\footnotesize Values are mean$\pm$SD
[minimum--maximum] unless otherwise stated. Demographic values use 23 NC and
51 MCI participants because no demographic row matched MCI 302. Physiological
recording values use all 23 NC and 52 MCI participants. $^{a}$MCI $n=50$ for
the geriatric depression score. Medication denominators reflect available
entries. BP, blood pressure; CBFV, cerebral blood flow velocity; CVRi, mean
MAP divided by mean CBFV; MMSE, Mini-Mental State Examination; MoCA, Montreal
Cognitive Assessment; SD, standard deviation.}\\
\end{longtable}

\subsection{Model dependence of gain}

The primary methodological result was that simultaneous modeling changed
pathway-specific gain, but the size and frequency pattern of that change were
not identical for MAP and $\COtwo$ (Table~\ref{tab:model_gain}). In the MAP
pathway, MISO gain was lower than SISO gain in VVLF for both groups. NC also
showed a small increase in MISO MAP gain in LF. The mean NC LF difference was
0.065, and the median subject-level MISO/SISO ratio was 1.028, indicating that
the statistically detectable effect was small in absolute terms.

The more consistent changes occurred in the $\COtwo$ pathway. NC MISO
$\COtwo$ gain was lower than SISO gain in LF and HF. MCI MISO $\COtwo$ gain
was lower in VLF, LF, and HF. The median MISO/SISO ratio in significant
$\COtwo$ comparisons ranged from 0.52 to 0.74. Thus, the separate $\COtwo$-only
model generally produced larger gain estimates than the conditional
$\COtwo$ estimate after MAP was included.

No MISO-versus-SISO phase comparison survived BH adjustment at the band level.
Partial coherence and ordinary coherence were also compared, but these values
were not treated as a test of model performance because they answer different
questions.

\begin{landscape}
\scriptsize
\begin{longtable}{p{0.48in}p{0.48in}p{0.48in}p{0.62in}p{0.62in}p{1.18in}p{0.62in}p{0.62in}}
\caption{Primary within-subject comparison of MISO and SISO gain.}
\label{tab:model_gain}\\
\toprule
Group & Path & Band & MISO & SISO & MISO$-$SISO (95\% CI) & Median ratio & BH $P$ \\
\midrule
\endfirsthead
\toprule
Group & Path & Band & MISO & SISO & MISO$-$SISO (95\% CI) & Median ratio & BH $P$ \\
\midrule
\endhead
NC & MAP & VVLF & 1.092 & 1.332 & $-0.240$ ($-0.384,-0.096$) & 0.837 & \textbf{0.0052} \\
NC & MAP & VLF  & 1.288 & 1.270 & $0.018$ ($-0.060,0.096$) & 0.997 & 0.6370 \\
NC & MAP & LF   & 2.111 & 2.046 & $0.065$ ($0.010,0.120$) & 1.028 & \textbf{0.0464} \\
NC & MAP & HF   & 2.343 & 2.331 & $0.012$ ($-0.036,0.060$) & 1.014 & 0.6370 \\
NC & $\COtwo$ & VVLF & 4.127 & 4.370 & $-0.243$ ($-0.675,0.189$) & 0.957 & 0.3149 \\
NC & $\COtwo$ & VLF  & 2.684 & 3.170 & $-0.486$ ($-1.087,0.116$) & 0.816 & 0.1444 \\
NC & $\COtwo$ & LF   & 1.761 & 2.498 & $-0.737$ ($-1.067,-0.407$) & 0.662 & \textbf{0.00052} \\
NC & $\COtwo$ & HF   & 2.695 & 4.944 & $-2.250$ ($-3.145,-1.354$) & 0.523 & \textbf{0.00017} \\
\addlinespace
MCI & MAP & VVLF & 0.837 & 1.020 & $-0.183$ ($-0.290,-0.076$) & 0.844 & \textbf{0.0038} \\
MCI & MAP & VLF  & 1.192 & 1.253 & $-0.061$ ($-0.116,-0.006$) & 0.971 & 0.0539 \\
MCI & MAP & LF   & 1.779 & 1.769 & $0.010$ ($-0.028,0.047$) & 0.994 & 0.6370 \\
MCI & MAP & HF   & 2.034 & 1.995 & $0.039$ ($-0.001,0.079$) & 1.015 & 0.0863 \\
MCI & $\COtwo$ & VVLF & 3.089 & 3.328 & $-0.239$ ($-0.493,0.014$) & 0.982 & 0.0931 \\
MCI & $\COtwo$ & VLF  & 2.066 & 2.850 & $-0.784$ ($-1.106,-0.462$) & 0.743 & \textbf{0.000083} \\
MCI & $\COtwo$ & LF   & 1.599 & 2.892 & $-1.293$ ($-1.718,-0.869$) & 0.620 & \textbf{0.0000021} \\
MCI & $\COtwo$ & HF   & 2.823 & 5.458 & $-2.635$ ($-4.275,-0.995$) & 0.682 & \textbf{0.0052} \\
\bottomrule
\multicolumn{8}{p{8.2in}}{\footnotesize Values are group means. Gain units are
percentage points of baseline CBFV per mmHg. Ratios were
calculated within subject as MISO/SISO and summarized by the median. BH
adjustment was performed across the 16 primary gain tests. Bold values are
BH-adjusted $P<0.05$. CI, confidence interval; HF, high frequency; LF, low
frequency; MCI, mild cognitive impairment; NC, normal control; VLF,
very-low frequency; VVLF, very-very-low frequency.}\\
\end{longtable}
\end{landscape}

% AUTHOR ACTION: Use the actual batch statistics output to make this figure.
\figureplaceholder[0.28\textheight]
{Paired subject-level MISO and SISO gain comparisons}
{Use four panels for MAP VVLF, MAP LF, $\COtwo$ LF, and $\COtwo$ HF, or split
MAP and $\COtwo$ into two figures. Plot each participant's SISO and MISO band
values connected by a line, with NC and MCI shown separately. Annotate the
mean paired difference, 95\% CI, median ratio, raw $P$, and BH-adjusted $P$.
The figure should show effect size and variability rather than only a
significance symbol.}
{fig:model_gain}

\subsection{Coherence structure and pathway balance}

Input coherence was modest across bands and similar between groups
(Table~\ref{tab:coherence}). Conditional $\COtwo$ partial coherence was
comparable to or greater than MAP partial coherence in VVLF. MAP partial
coherence became larger than $\COtwo$ partial coherence from VLF through HF.
This frequency pattern was present in both NC and MCI. Multiple coherence was
between approximately 0.58 and 0.75 across bands, but it was interpreted as
in-sample representation by both inputs rather than predictive validation.

The MAP-minus-$\COtwo$ partial-coherence balance was negative in VVLF and
positive in VLF, LF, and HF. The balance did not differ between NC and MCI in
any band after BH adjustment. These findings indicate a similar broad
frequency-dependent balance of conditional associations in both groups.

\begin{table}[H]
\centering
\small
\caption{Descriptive band coherence values.}
\label{tab:coherence}
\begin{tabular}{llcccc}
\toprule
Metric & Group & VVLF & VLF & LF & HF \\
\midrule
Input MAP--$\COtwo$ coherence & NC  & 0.314 & 0.250 & 0.216 & 0.209 \\
                              & MCI & 0.262 & 0.264 & 0.215 & 0.190 \\
MISO MAP partial coherence    & NC  & 0.335 & 0.560 & 0.592 & 0.656 \\
                              & MCI & 0.276 & 0.548 & 0.671 & 0.656 \\
MISO $\COtwo$ partial coherence & NC  & 0.409 & 0.346 & 0.221 & 0.225 \\
                                & MCI & 0.385 & 0.268 & 0.233 & 0.243 \\
MISO multiple coherence       & NC  & 0.645 & 0.711 & 0.681 & 0.733 \\
                              & MCI & 0.582 & 0.695 & 0.746 & 0.731 \\
SISO MAP ordinary coherence   & NC  & 0.388 & 0.542 & 0.584 & 0.647 \\
                              & MCI & 0.308 & 0.569 & 0.665 & 0.644 \\
SISO $\COtwo$ ordinary coherence & NC  & 0.456 & 0.300 & 0.198 & 0.209 \\
                                  & MCI & 0.412 & 0.286 & 0.210 & 0.193 \\
MAP minus $\COtwo$ partial coherence & NC  & $-0.074$ & 0.214 & 0.370 & 0.431 \\
                                      & MCI & $-0.109$ & 0.281 & 0.438 & 0.413 \\
\bottomrule
\end{tabular}
\end{table}

% AUTHOR ACTION: The ordinary coherence line belongs only on SISO panels.
\figureplaceholder[0.27\textheight]
{Group coherence structure across frequency}
{For each group, show MAP--$\COtwo$ input coherence, MAP partial coherence,
$\COtwo$ partial coherence, and multiple coherence as frequency curves with
group mean and between-subject SD. In separate SISO panels, show ordinary MAP
and $\COtwo$ coherence with the CARNet reference line. Do not draw the
ordinary-coherence threshold on partial or multiple coherence.}
{fig:coherence}

\subsection{NC versus MCI comparisons}

No prespecified band-level NC--MCI comparison of MISO gain or partial coherence
survived BH correction (Table~\ref{tab:group_primary}). The largest unadjusted
gain difference was in the MAP LF band, where NC had higher MISO gain than MCI
(difference 0.332, 95\% CI 0.024 to 0.640, Hedges' $g=0.622$, raw $P=0.0355$).
After adjustment across the primary family, the BH-adjusted $P$ value was
0.335. The largest partial-coherence difference was in the $\COtwo$ VLF band
(NC minus MCI 0.078, 95\% CI 0.003 to 0.153, raw $P=0.0418$,
BH-adjusted $P=0.335$). No band-level MISO phase comparison survived the
secondary phase adjustment.

\begin{landscape}
\scriptsize
\begin{longtable}{p{0.48in}p{0.60in}p{0.45in}p{0.62in}p{0.62in}p{1.20in}p{0.54in}p{0.58in}}
\caption{Primary NC versus MCI comparisons of MISO gain and partial coherence.}
\label{tab:group_primary}\\
\toprule
Path & Metric & Band & NC & MCI & NC$-$MCI (95\% CI) & Hedges $g$ & BH $P$ \\
\midrule
\endfirsthead
\toprule
Path & Metric & Band & NC & MCI & NC$-$MCI (95\% CI) & Hedges $g$ & BH $P$ \\
\midrule
\endhead
MAP & Gain & VVLF & 1.092 & 0.837 & 0.255 ($-0.140,0.649$) & 0.344 & 0.469 \\
MAP & Gain & VLF  & 1.288 & 1.192 & 0.097 ($-0.092,0.285$) & 0.229 & 0.495 \\
MAP & Gain & LF   & 2.111 & 1.779 & 0.332 ($0.024,0.640$) & 0.622 & 0.335 \\
MAP & Gain & HF   & 2.343 & 2.034 & 0.309 ($-0.071,0.689$) & 0.477 & 0.344 \\
MAP & Partial coh. & VVLF & 0.335 & 0.276 & 0.059 ($-0.034,0.153$) & 0.330 & 0.469 \\
MAP & Partial coh. & VLF  & 0.560 & 0.548 & 0.012 ($-0.089,0.112$) & 0.061 & 0.870 \\
MAP & Partial coh. & LF   & 0.592 & 0.671 & $-0.079$ ($-0.174,0.015$) & $-0.414$ & 0.344 \\
MAP & Partial coh. & HF   & 0.656 & 0.656 & $-0.0002$ ($-0.086,0.086$) & $-0.001$ & 0.996 \\
\addlinespace
$\COtwo$ & Gain & VVLF & 4.127 & 3.089 & 1.038 ($-0.214,2.291$) & 0.493 & 0.344 \\
$\COtwo$ & Gain & VLF  & 2.684 & 2.066 & 0.619 ($-0.426,1.663$) & 0.380 & 0.470 \\
$\COtwo$ & Gain & LF   & 1.761 & 1.599 & 0.163 ($-0.291,0.616$) & 0.186 & 0.631 \\
$\COtwo$ & Gain & HF   & 2.695 & 2.823 & $-0.128$ ($-1.095,0.839$) & $-0.064$ & 0.870 \\
$\COtwo$ & Partial coh. & VVLF & 0.409 & 0.385 & 0.024 ($-0.072,0.121$) & 0.112 & 0.759 \\
$\COtwo$ & Partial coh. & VLF  & 0.346 & 0.268 & 0.078 ($0.003,0.153$) & 0.548 & 0.335 \\
$\COtwo$ & Partial coh. & LF   & 0.221 & 0.233 & $-0.011$ ($-0.042,0.019$) & $-0.184$ & 0.631 \\
$\COtwo$ & Partial coh. & HF   & 0.225 & 0.243 & $-0.017$ ($-0.049,0.014$) & $-0.281$ & 0.487 \\
\bottomrule
\multicolumn{8}{p{8.2in}}{\footnotesize BH adjustment was performed across
the 16 primary group tests. Gain units are percentage points of baseline CBFV
per mmHg; coherence is unitless. No comparison met BH-adjusted $P<0.05$. CI,
confidence interval; coh., coherence; HF, high frequency; LF, low frequency;
MCI, mild cognitive impairment; NC, normal control; VLF, very-low frequency;
VVLF, very-very-low frequency.}\\
\end{longtable}
\end{landscape}

% AUTHOR ACTION: This should be a distribution figure, not only mean bars.
\figureplaceholder[0.25\textheight]
{Individual MISO values used in the NC--MCI comparisons}
{Plot individual NC and MCI values for the most relevant prespecified outcomes:
MAP gain and partial coherence in VVLF/VLF/LF and $\COtwo$ gain and partial
coherence in VVLF/VLF/LF. Use a central estimate with a 95\% CI, not mean
$\pm$SD error bars that extend beyond bounded coherence limits. Report raw and
BH-adjusted $P$ values. A supplementary figure can contain all four bands.}
{fig:group_comparison}

\subsection{Exploratory frequency-wise comparisons}

The frequency-wise NC--MCI analysis identified one bin that survived BH
adjustment within its comparison curve. At 0.1484375~Hz, MISO MAP gain was
higher in NC than MCI (NC 2.462, MCI 1.710; difference 0.752, 95\% CI 0.334 to
1.170; Hedges' $g=0.959$; raw $P=0.000813$; BH-adjusted $P=0.0366$). Several
neighboring bins had raw $P<0.05$, but only the 0.1484-Hz bin survived the
curve-wise correction. No frequency-wise group phase or partial-coherence bin
survived its curve-wise BH correction.

Frequency-wise model comparisons showed broad MISO--SISO differences in
$\COtwo$ gain, especially in MCI, and more limited differences in MAP gain.
These curves were consistent with the band results: simultaneous modeling had
a stronger effect on $\COtwo$ gain than on MAP gain. No frequency-wise phase
comparison survived its curve-wise correction.

% AUTHOR ACTION: This is the main frequency-wise inferential figure.
\figureplaceholder[0.32\textheight]
{Exploratory frequency-wise NC--MCI comparison of MISO MAP gain}
{Top panel: NC and MCI group mean MISO MAP gain with between-subject SD or a
clearly labeled 95\% CI. Middle panel: NC-minus-MCI mean difference with a
zero line. Bottom panel: raw and BH-adjusted $P$ curves on a $-\log_{10}(P)$
axis so small values are visible, with a horizontal line at
$-\log_{10}(0.05)$. Mark 0.1484~Hz and draw the prespecified band boundaries.
State in the caption that BH correction was within this 45-bin curve and that
the result is exploratory.}
{fig:frequency_group}

% AUTHOR ACTION: Consider supplement if the main paper becomes figure-heavy.
\figureplaceholder[0.32\textheight]
{Frequency-wise MISO versus SISO gain comparisons}
{Create separate NC and MCI figures or a compact multi-panel figure. For MAP
and $\COtwo$, show the MISO and SISO group curves, their paired difference,
and the raw/BH-adjusted $P$ curves. Use gain, not phase, in the main version
because no phase comparison survived adjustment.}
{fig:frequency_model}

\subsection{Input coherence and model-related gain changes}

MAP--$\COtwo$ input coherence was positively associated with the absolute
MISO-minus-SISO gain change in several exploratory comparisons. In MCI, the
associations included MAP VVLF ($\rho=0.551$, BH-adjusted $P=0.00038$), MAP LF
($\rho=0.363$, BH-adjusted $P=0.0338$), and $\COtwo$ VVLF ($\rho=0.365$,
BH-adjusted $P=0.0338$). In the pooled cohort, significant associations were
present for MAP VVLF ($\rho=0.537$, BH-adjusted $P=2.46\times10^{-5}$), MAP
VLF ($\rho=0.293$, BH-adjusted $P=0.0379$), MAP LF ($\rho=0.283$,
BH-adjusted $P=0.0421$), $\COtwo$ VVLF ($\rho=0.375$, BH-adjusted $P=0.0080$),
and $\COtwo$ VLF ($\rho=0.347$, BH-adjusted $P=0.0145$).

These results support the idea that the two models differ more when the inputs
share more frequency content. However, they were not treated as independent
biological validation because both input coherence and gain change are
calculated from the same underlying spectral estimates.

\figureplaceholder[0.25\textheight]
{Input coherence and the size of model-related gain changes}
{For all four bands, plot each subject's band-averaged MAP--$\COtwo$ input
coherence against the absolute MISO-minus-SISO gain change. Use separate MAP
and $\COtwo$ pathway panels, display NC and MCI with different colors, and
annotate $n$, Spearman $\rho$, raw $P$, and BH-adjusted $P$. Keep the word
``exploratory'' in the title or caption.}
{fig:input_association}

\subsection{Numerical conditioning}

The raw input spectral matrix was more poorly conditioned at some high-frequency
points. The median HF condition number was 10.9 in NC and 15.4 in MCI. Values
greater than 100 occurred in 60 of 437 NC and 126 of 988 MCI HF
frequency-by-subject observations. The maximum values were approximately
7,467 in NC and 17,366 in MCI. High values tended to occur when $\PETCOtwo$
power was low or MAP and $\PETCOtwo$ powers were highly imbalanced.

Because the raw condition number depends on the units and scaling of both
inputs, these observations do not establish severe physiological collinearity.
They do show that some high-frequency MISO estimates may be sensitive to small
spectral changes. No regularization or condition-number exclusion was applied.

% AUTHOR ACTION: This sensitivity analysis should be resolved with the PI.
\draftnote{Before submission, decide whether to add a scale-invariant
conditioning diagnostic based on the normalized input spectral matrix and
whether a prespecified ridge sensitivity analysis is needed. Do not add an
arbitrary ridge constant after seeing the results. If regularization is used,
justify its scale and show that the main conclusions are stable.}

\figureplaceholder[0.23\textheight]
{Input spectral-matrix conditioning diagnostics}
{Suggested supplementary figure: condition number versus frequency for each
group, plus condition number versus $\PETCOtwo$ power and MAP/$\PETCOtwo$
power ratio. Clearly state that the plotted raw condition number is
scale-dependent.}
{fig:conditioning}

% ============================================================================
\section{Discussion}
% ============================================================================

\subsection{Principal findings}

This study asked two related but distinct questions. First, does simultaneous
MAP and $\PETCOtwo$ modeling materially change pathway-specific transfer
function estimates? Second, after the pathways are modeled together, do the
conditional estimates differ between NC and MCI? The answer to the first
question was yes. The answer to the second was not at the prespecified band
level.

The clearest methodological result was the reduction in MISO $\COtwo$ gain
relative to $\COtwo$-only SISO gain. This reduction was present in NC at LF and
HF and in MCI at VLF, LF, and HF. MISO MAP gain also differed from SISO MAP
gain, but the changes were smaller and concentrated in VVLF. These findings
show that the one-input and two-input models are not interchangeable. They do
not show that the MISO model is universally superior. Rather, the models
partition shared input--output information differently.

The main biological hypothesis was not supported. None of the prespecified
band-level NC--MCI differences in MISO gain, partial coherence, phase, or
pathway balance survived BH correction. This negative result is important.
It indicates that accounting for spontaneous $\PETCOtwo$ fluctuations did not
reveal a broad conditional MAP-pathway deficit in the present MCI cohort. The
finding is compatible with prior reports and systematic reviews showing mixed
or absent group differences in spontaneous dCA metrics
\cite{Gommer2012,Tarumi2014,Heutz2023}.

\subsection{Why SISO and MISO gain differed}

A SISO transfer function assigns all linear input--output cross-spectral
information to one input. If MAP and $\PETCOtwo$ are coupled, and both are
associated with CBFV, the separate SISO estimates may include shared
information. The MISO solution uses the full input spectral matrix and
estimates both transfer functions at the same time. The conditional coefficient
for one input can therefore decrease when variation shared with the other
input is accounted for.

The strong reduction in $\COtwo$ gain is consistent with this mechanism. The
result does not mean that $\COtwo$ is physiologically unimportant. Carbon
dioxide remains a potent regulator of cerebral blood flow
\cite{Hoiland2019}. It means that the $\COtwo$-only SISO gain was larger than
the coefficient assigned specifically to $\PETCOtwo$ after simultaneous MAP
modeling. In other words, part of the pairwise $\PETCOtwo$--CBFV relationship
was shared with MAP in the two-input representation.

The pooled positive associations between input coherence and absolute gain
change further support this interpretation. They are strongest in low
frequency bands, where prior multivariable work has also shown that gas-related
inputs can change pressure--flow estimates \cite{Panerai2000,Peng2008}.
Nevertheless, these associations should remain secondary because they are
partly mathematically coupled to the same spectral matrix used to calculate
the transfer functions.

\subsection{Frequency-dependent pathway balance}

The descriptive partial-coherence pattern was physiologically plausible but
should not be overinterpreted. Conditional $\COtwo$ coherence was comparable
to or greater than conditional MAP coherence in VVLF. From VLF through HF,
conditional MAP coherence was larger. This pattern suggests that slow CBFV
variation contains an important $\COtwo$-associated component, while MAP is
more strongly associated with CBFV at higher frequencies in the analyzed
range.

The pathway-balance result was similar in NC and MCI. It therefore provides a
useful description of frequency-dependent association but not evidence of an
MCI-specific shift. Moreover, partial coherence is not a direct measure of
causal contribution. Unmeasured inputs, feedback, measurement noise, and
nonlinear physiology can all affect the observed spectral relationships.

\subsection{Interpretation of the NC--MCI results}

The absence of BH-significant band differences should not be described as
proof that NC and MCI are the same. The confidence intervals show the range of
effects compatible with the current data. For example, NC MAP LF gain was
higher than MCI before multiplicity adjustment and had a moderate Hedges'
$g$. The adjusted analysis, however, did not provide enough evidence to
declare that result across the prespecified family. The correct interpretation
is that the study did not establish a robust band-level group difference.

The isolated 0.1484-Hz MAP-gain difference deserves follow-up but not a primary
claim. It survived BH adjustment across the bins of its own curve and was
surrounded by several smaller raw $P$ values, which is more credible than a
completely isolated numerical spike. Still, the frequency-wise analyses were
exploratory, multiple metric/pathway curves were examined, and the finding was
not a prespecified single-frequency hypothesis. It should therefore be used
to guide a future targeted analysis or replication rather than to redefine
the main conclusion of this study.

\subsection{Coherence as a benchmark rather than a filter}

The present analysis retained all frequency bins and used ordinary SISO
coherence as a visual benchmark. This choice separates data description from
data deletion. With three Welch windows, the CARNet ordinary-coherence
reference was 0.51 \cite{Panerai2023}. Much of the $\COtwo$ ordinary coherence,
especially outside VVLF, was below this reference. Consequently, the high-
frequency $\COtwo$ gain and phase estimates should be interpreted cautiously
even when MISO and SISO gains differed statistically.

The ordinary threshold does not transfer directly to partial or multiple
coherence. Applying the same number to those metrics would imply a sampling
distribution that has not been established here. Showing the ordinary line
only on SISO coherence plots is therefore more defensible than filtering the
full analysis or drawing the same line on every coherence statistic.

\subsection{Phase findings}

No phase comparison survived multiplicity correction. This is not surprising
given the circular variability of phase, the limited number of Welch windows,
and abrupt branch changes near $-\pi$ and $\pi$. Circular averaging prevents
the specific error of treating those boundaries as far apart, and testing the
wrapped values avoids making inferential results depend on the chosen display
branch. Unwrapped curves remain useful for visualization, but they should be
shown as the unwrapped circular mean with the same circular standard deviation.

\subsection{Multiple coherence and prediction}

Multiple coherence was generally higher than either single pairwise coherence
because two inputs were fitted to the same output on the same data. This is
expected and does not by itself demonstrate that MISO predicts CBFV better in
new data. A meaningful prediction comparison would estimate the spectra or
model on a training segment and evaluate reconstructed CBFV on a separate
test segment. Most recordings in this cohort provided only three 128-second
Welch windows, leaving little independent data for both stable estimation and
testing.

Cross-validation is therefore better suited to a future supplemental analysis
using longer recordings. A 20-second holdout after the 256-second minimum
would contain relatively little low-frequency information and would not
provide a strong test of VVLF or VLF prediction. Longer recordings or an
experimental design with repeated baseline epochs would allow a more
defensible training/test separation.

\subsection{Numerical stability}

The MISO solution was obtained without regularization. This keeps the current
analysis transparent, but some raw spectral matrices were poorly conditioned,
especially at high frequencies with low $\PETCOtwo$ power. The largest gain
changes also occurred in higher-frequency $\COtwo$ bands, so numerical
stability and physiological interpretation cannot be separated completely.

A future sensitivity analysis should first normalize the input spectral
matrix so that conditioning is not dominated by the different units of MAP
and $\PETCOtwo$. If regularization is then needed, the ridge scale should be
defined before viewing group results and justified relative to the normalized
matrix or measurement-noise level. The smallest ridge that produces stable
solutions across a prespecified sensitivity range would be preferable to an
arbitrary constant. The current paper should not claim that such a sensitivity
analysis was performed unless it is actually added.

\subsection{Strengths}

This study has several strengths. First, it directly compares SISO and MISO
estimates within the same participants, so model-related changes are not
confounded by different cohorts. Second, MAP, $\PETCOtwo$, and CBFV share one
spectral estimation step, and identical smoothing is applied to all auto- and
cross-spectra. Third, phase means and dispersion use circular rather than
linear statistics. Fourth, inferential tests use one band value per subject
and do not treat frequency bins as independent subjects. Fifth, raw and
BH-adjusted $P$ values, confidence intervals, standardized effects, and median
gain ratios are exported with the subject-level values used in each test.
Finally, the workflow is parameterized from one main file and includes tests
for its major numerical and export components.

\subsection{Limitations}

Several limitations affect interpretation. The study was observational, so
the transfer functions describe association and do not establish causal MAP-
or $\COtwo$-mediated effects. $\PETCOtwo$ is a surrogate for arterial carbon
dioxide and may be influenced by ventilation, dead space, and measurement
delay. Transcranial Doppler measures velocity rather than volumetric flow and
assumes that the insonated arterial diameter is sufficiently stable
\cite{Aaslid1982}. Only baseline spontaneous fluctuations were analyzed, which
may provide limited input power compared with repeated squat--stand or
controlled $\COtwo$ challenges.

Most recordings produced only three Welch windows. This meets the configured
minimum but limits spectral precision, raises the ordinary-coherence reference,
and prevents strong cross-validation. The $\PETCOtwo$ startup rule was designed
for the available exports and should be checked in new datasets. The MISO
model included two inputs, but cerebral blood flow is affected by other
variables, including oxygen, cardiac output, autonomic activity, metabolism,
and neurovascular coupling. Feedback between MAP, ventilation, and the
cerebral circulation also challenges a simple feed-forward linear model.

The raw input-matrix condition number was scale-dependent, and no ridge
sensitivity analysis was performed. Frequency-domain smoothing makes
neighboring bins more similar, so frequency-wise $P$ values should not be read
as independent discoveries. Finally, the NC and MCI sample sizes were unequal,
one MCI participant lacked demographic data, and the diagnostic composition
of MCI must be described more fully from the parent protocol.

\subsection{Future work}

The next study should prioritize longer recordings or repeated baseline
epochs, prospectively defined model-stability criteria, and out-of-sample
prediction. Controlled MAP oscillations and controlled $\COtwo$ challenges
could provide stronger and more separable inputs than spontaneous baseline
variation. A normalized conditioning analysis and prespecified regularization
sensitivity test would clarify the stability of high-frequency coefficients.
The 0.1484-Hz MAP-gain finding should be tested as a targeted hypothesis in an
independent cohort rather than searched for again across the same frequency
range.

Future biological analyses could also test associations with continuous
cognitive scores, MCI subtype, vascular risk, medication exposure, or imaging
markers rather than relying only on a binary NC--MCI comparison. Those analyses
should be defined before testing and should account for age, sex, education,
and other relevant covariates when the sample size supports them.

% ============================================================================
\section{Conclusion}
% ============================================================================

Simultaneously modeling MAP and $\PETCOtwo$ materially changed transfer
function gain estimates, with the largest and most consistent changes in the
$\COtwo$ pathway. This supports the methodological argument that separate
MAP-only and $\COtwo$-only SISO models can combine information that is
partitioned differently in a two-input model. In contrast, the current data
did not show a robust prespecified band-level difference in conditional MAP-
or $\COtwo$-associated CBFV dynamics between NC and MCI. The results therefore
support pathway separation as a useful analysis framework, while they do not
support a broad MCI-specific impairment under resting spontaneous conditions.
The exploratory difference at 0.1484~Hz and the high-frequency conditioning
findings provide clear directions for replication and sensitivity analysis.

% ============================================================================
\section*{Declarations}
% ============================================================================

\subsection*{Ethics approval and consent to participate}
% AUTHOR ACTION: Insert PI-approved ethics statement.
\draftnote{Insert the institutional review board name and approval number,
and state how written informed consent was obtained.}

\subsection*{Data availability}
% AUTHOR ACTION: Define what can be shared and under what agreement.
\draftnote{State whether de-identified physiological data can be shared,
whether access requires a data-use agreement, and where the synthetic example
data and analysis code will be deposited.}

\subsection*{Code availability}
% AUTHOR ACTION: Add final URL and release.
The MATLAB analysis workflow will be made available at [repository URL and
archived release to be added].

\subsection*{Funding}
% AUTHOR ACTION: Insert grant and institutional support.
[Funding information to be added.]

\subsection*{Conflicts of interest}
% AUTHOR ACTION: Confirm with every author.
[Conflict-of-interest statement to be added.]

\subsection*{Author contributions}
% AUTHOR ACTION: Use CRediT roles if required by the target journal.
[Author contribution statement to be added.]

\subsection*{Acknowledgments}
[Acknowledgments to be added.]

% ============================================================================
\appendix
\section{Recommended figure set and placement}
% ============================================================================

The following set keeps the main paper focused while preserving the statistical
evidence:

\begin{enumerate}[leftmargin=0.45in]
    \item \textbf{Main Figure 1:} Conceptual SISO-versus-MISO schematic
    (Figure~\ref{fig:concept}).
    \item \textbf{Main Figure 2:} Paired subject-level MISO-versus-SISO gain
    comparisons (Figure~\ref{fig:model_gain}).
    \item \textbf{Main Figure 3:} Coherence structure across frequency
    (Figure~\ref{fig:coherence}).
    \item \textbf{Main Figure 4:} Individual NC and MCI MISO band values
    (Figure~\ref{fig:group_comparison}).
    \item \textbf{Main or supplemental Figure 5:} Exploratory 0.1484-Hz
    MAP-gain result with adjusted $P$ curve (Figure~\ref{fig:frequency_group}).
    \item \textbf{Supplemental Figure S1:} Full frequency-wise MISO-versus-SISO
    gain results (Figure~\ref{fig:frequency_model}).
    \item \textbf{Supplemental Figure S2:} Input coherence associations
    (Figure~\ref{fig:input_association}).
    \item \textbf{Supplemental Figure S3:} Conditioning diagnostics
    (Figure~\ref{fig:conditioning}).
\end{enumerate}

Representative single-subject raw signals and subject-specific transfer
function curves could be added as a supplemental quality-control figure. They
are useful for showing what the input data and model output look like, but they
should not replace group distributions or confidence intervals.

\section{Items that must be resolved before submission}

\begin{enumerate}[leftmargin=0.45in]
    \item Confirm participant diagnosis, recruitment, exclusions, IRB, consent,
    acquisition devices, posture, and recording protocol from the parent study.
    \item Decide whether the 0-Hz bin remains in exploratory frequency-wise
    tests.
    \item Decide with the PI whether normalized conditioning and a prespecified
    ridge sensitivity analysis are needed. The main reported solve is currently
    unregularized.
    \item Decide whether the primary family should remain all four bands or
    emphasize VVLF/VLF based on the biological hypothesis. This must be decided
    prospectively for any rerun and reported transparently.
    \item Confirm whether the Geriatric Depression Scale belongs in the main
    participant table and whether medication variables are sufficiently
    complete for publication.
    \item Confirm how to describe MCI 302, which has TFA results but no matching
    row in the demographic workbook.
    \item Add a permanent code release and data-availability statement.
    \item Select the target journal and reformat title page, abstract length,
    headings, tables, references, and figure limits.
\end{enumerate}

% ============================================================================
% REFERENCES
% ============================================================================

\begin{thebibliography}{99}

\bibitem{Claassen2021}
Claassen JAHR, Thijssen DHJ, Panerai RB, Faraci FM.
Regulation of cerebral blood flow in humans: physiology and clinical
implications of autoregulation.
\textit{Physiol Rev}. 2021;101(4):1487--1559.
doi:10.1152/physrev.00022.2020.

\bibitem{Panerai2023}
Panerai RB, Brassard P, Burma JS, et al.; Cerebrovascular Research Network.
Transfer function analysis of dynamic cerebral autoregulation: A CARNet white
paper 2022 update.
\textit{J Cereb Blood Flow Metab}. 2023;43(1):3--25.
doi:10.1177/0271678X221119760.

\bibitem{Kisler2017}
Kisler K, Nelson AR, Montagne A, Zlokovic BV.
Cerebral blood flow regulation and neurovascular dysfunction in Alzheimer
disease.
\textit{Nat Rev Neurosci}. 2017;18(7):419--434.
doi:10.1038/nrn.2017.48.

\bibitem{Beishon2017}
Beishon L, Haunton VJ, Panerai RB, Robinson TG.
Cerebral hemodynamics in mild cognitive impairment: A systematic review.
\textit{J Alzheimers Dis}. 2017;59(1):369--385.
doi:10.3233/JAD-170181.

\bibitem{Gommer2012}
Gommer ED, Martens EGHJ, Aalten P, et al.
Dynamic cerebral autoregulation in subjects with Alzheimer's disease, mild
cognitive impairment, and controls: evidence for increased peripheral
vascular resistance with possible predictive value.
\textit{J Alzheimers Dis}. 2012;30(4):805--813.
doi:10.3233/JAD-2012-111628.

\bibitem{Tarumi2014}
Tarumi T, Dunsky DI, Khan MA, et al.
Dynamic cerebral autoregulation and tissue oxygenation in amnestic mild
cognitive impairment.
\textit{J Alzheimers Dis}. 2014;41(3):765--778.
doi:10.3233/JAD-132018.

\bibitem{Heutz2023}
Heutz R, Claassen J, Feiner S, et al.
Dynamic cerebral autoregulation in Alzheimer's disease and mild cognitive
impairment: A systematic review.
\textit{J Cereb Blood Flow Metab}. 2023;43(8):1223--1236.
doi:10.1177/0271678X231173449.

\bibitem{Hoiland2019}
Hoiland RL, Fisher JA, Ainslie PN.
Regulation of the cerebral circulation by arterial carbon dioxide.
\textit{Compr Physiol}. 2019;9(3):1101--1154.
doi:10.1002/cphy.c180021.

\bibitem{Panerai2000}
Panerai RB, Simpson DM, Deverson ST, Mahony P, Hayes P, Evans DH.
Multivariate dynamic analysis of cerebral blood flow regulation in humans.
\textit{IEEE Trans Biomed Eng}. 2000;47(3):419--423.
doi:10.1109/10.827312.

\bibitem{Peng2008}
Peng T, Rowley AB, Ainslie PN, Poulin MJ, Payne SJ.
Multivariate system identification for cerebral autoregulation.
\textit{Ann Biomed Eng}. 2008;36(2):308--320.
doi:10.1007/s10439-007-9412-9.

\bibitem{Welch1967}
Welch PD.
The use of fast Fourier transform for the estimation of power spectra: A
method based on time averaging over short, modified periodograms.
\textit{IEEE Trans Audio Electroacoust}. 1967;15(2):70--73.
doi:10.1109/TAU.1967.1161901.

\bibitem{MardiaJupp2000}
Mardia KV, Jupp PE.
\textit{Directional Statistics}.
Chichester, UK: Wiley; 2000.

\bibitem{Berens2009}
Berens P.
CircStat: A MATLAB toolbox for circular statistics.
\textit{J Stat Softw}. 2009;31(10):1--21.
doi:10.18637/jss.v031.i10.

\bibitem{Welch1947}
Welch BL.
The generalization of Student's problem when several different population
variances are involved.
\textit{Biometrika}. 1947;34(1--2):28--35.
doi:10.1093/biomet/34.1-2.28.

\bibitem{Hedges1981}
Hedges LV.
Distribution theory for Glass's estimator of effect size and related
estimators.
\textit{J Educ Stat}. 1981;6(2):107--128.
doi:10.2307/1164588.

\bibitem{BenjaminiHochberg1995}
Benjamini Y, Hochberg Y.
Controlling the false discovery rate: A practical and powerful approach to
multiple testing.
\textit{J R Stat Soc Series B Stat Methodol}. 1995;57(1):289--300.
doi:10.1111/j.2517-6161.1995.tb02031.x.

\bibitem{Aaslid1982}
Aaslid R, Markwalder TM, Nornes H.
Noninvasive transcranial Doppler ultrasound recording of flow velocity in
basal cerebral arteries.
\textit{J Neurosurg}. 1982;57(6):769--774.
doi:10.3171/jns.1982.57.6.0769.

\bibitem{Folstein1975}
Folstein MF, Folstein SE, McHugh PR.
``Mini-mental state'': A practical method for grading the cognitive state of
patients for the clinician.
\textit{J Psychiatr Res}. 1975;12(3):189--198.
doi:10.1016/0022-3956(75)90026-6.

\bibitem{Nasreddine2005}
Nasreddine ZS, Phillips NA, B\'edirian V, et al.
The Montreal Cognitive Assessment, MoCA: A brief screening tool for mild
cognitive impairment.
\textit{J Am Geriatr Soc}. 2005;53(4):695--699.
doi:10.1111/j.1532-5415.2005.53221.x.

\end{thebibliography}

\end{document}
